\typesetchapter{Chapter 9}{Reggane}

\begin{room}

--- There is a glass in the long wall of this room, Paul, and a smaller room behind it, and I have sat in the smaller room morning after morning and watched you take him apart with a patience that could have made me weep. He has given you the bowl and the pendulum and the girl and the drowned, and the raising into Light that I asked to be allowed to see and saw, through the glass, and wept at, where no one could tell me not to. And I have given him a name across a tea tray and nothing else. We are coming, the two of us, to the same desert from opposite sides of it, and I will not sit behind glass and let him describe the crossing to a machine from one side only. So I have come round to the door. I am going to give my half of it, into your machine, in his hearing, from the beginning, the way he gave his. He may stop me wherever he likes.

--- \ldots{} You are the woman who pays for the tea.

--- I am the reason there is a house to drink it in, and a machine to catch you, and a man patient enough to believe you. You have told the whole of your life to my employee, day after day; it is the barest courtesy that you should now have mine. And it is more than courtesy: the record is worth nothing unless both hands that touched the thing are set down in it. Ask me whatever you like when I am done.

--- \ldots{} Then come out from behind your glass, and sit where I can see you. Since the morning you walked through this room and named us your two clevernesses, I have wanted to know the shape of the hand that reached into my mountain. Begin where it begins.

\end{room}

I was twenty-seven the year the French lit their first bomb in the Sahara, and I was one of perhaps forty people in the world positioned to hear it properly, and I was the only one of the forty who was bored by it.\par

Understand what a spectacle it was. A country that had lost a war and an empire and most of its dignity within living memory had decided to buy all three back in a single morning with a tower of steel and a sun it made itself, in a corner of a desert it still, for a few more years, called France. Reggane. They put the device on a hundred metres of scaffold, and they filled the ground around it with instruments, and they invited the whole nervous world to listen, and the whole nervous world listened, because a nation that can make its own sun at will is a nation you keep an ear on afterward. Gerboise blue, they called it. A blue desert rat. The French have always named their monsters as if for a nursery.\par

My husband put me there. He was a physicist of the kind the new French state was collecting like uranium ore in those years --- rare, refined at great cost, and handled by men who did not understand him and could not do without him. He believed in the bomb the way other men believe in God: as a thing too large to be moral, and therefore beneath morality, and therefore his to serve without stain. He was clever and vain and, in the one narrow chapel of his cleverness, close to great. I married him for the chapel. I have never pretended otherwise, least of all to him.\par

What he had that I wanted was a way of thinking I did not yet own. He thought in chains. A single splitting atom is nothing; it would not warm your hands. But it throws two more, and each of those two throws two, and if you have gathered enough of the right stuff close enough together, the nothing becomes, in the space of a breath, the end of a city. He could hold that in his head as a plain fact. The largest thing was only the smallest thing, permitted to answer itself. I took that from him the way you take a tool from a bench when the owner has turned away, and I never gave it back, and it was the most valuable thing I ever stole, and I stole a great deal.\par

My father had given me the other half without knowing it. Edward Blackwood spent his life moving the buried wealth of the world from where it lay to where it would sell, and the deep Sahara was one of his cellars. Oil in the north. Iron. And, once the French wanted their own sun, the yellow rock, which he had been quietly gathering the rights to for a decade before anyone official admitted wanting it. My father was not a scientist. He was something rarer and, in his trade, more useful: a man with a nose. And for the last stretch of his working life his nose had been turned, without profit and against all his instincts, toward an empty quarter of that desert where nothing he touched ever came to anything. Concessions he bought and never worked. Surveys that came back wrong or came back not at all. Two expeditions he funded that turned around short of where he sent them, for reasons the men who led them could never afterward say to his face. He knew there was something out there. He was certain of it the way he was certain of a seam of ore under a hillside, by a pressure behind the sternum that had made him rich. And he chased it for thirty years and a fortune and died having never laid a hand on it. I loved him. I do not think he would have known what to do with it if he had found it. I did.\par

So. Between them, the two men who between them made me, I had the chain and the cellar. My husband taught me that the ground could carry the shape of a thing that had happened in it, out and out to any patient enough ear, and that the shape could be read for its cause. That was new then. Men were only just learning to tell a bomb from an earthquake by the drawing it left on a drum --- the bomb sharp at the front, all of its violence arriving at once, the way a lie arrives; the earthquake ragged and building and easing, the way the truth does. A whole science was being born in those years for one small dishonest purpose: so that the great powers could catch each other testing in secret. I was inside that science because my husband was, and I was better at the part of it he found tedious, which was the reading. He wanted to make the sun. I wanted to sit in the dark afterward with the drum-paper and learn what the desert had been keeping in its notes.\par

There is a grammar to that reading, and I had it cold before I was thirty. A natural quake announces itself by its untidiness: it builds, it wanders, it arrives in a long ragged crowd of waves that have taken different roads through the rock and come in out of step with one another, and it never once repeats itself, because the earth keeps no discipline. An explosion is the opposite kind of lie: it is too tidy, all of its force flung out in a single hard shove in the first instant and nothing after it, a shout with no sentence to follow. Between those two poles a trained reader learns to place almost everything the ground can do --- a cliff falling, a mine caving on its own galleries, a heavy truck on a bad road, a door slammed somewhere in the recorder's own building. The one thing a trained reader is never made ready for, because nothing in the honest earth produces it, is a figure that is both tidy and repeated: a small deliberate hand, struck and struck again across the years without one hair of variation. That is something standing on the earth and using it.\par

\begin{room}

--- When was this.

--- The first bomb? A lifetime ago. Yours. Do the arithmetic later; it is not the year that matters. It is that the ear existed, and was turned on your desert, and you did not know an ear could exist. Your whole peace rested on a thing that had stopped being true before you were born.

\end{room}

The bomb was on the tower and the world was holding its breath, and I was in a tent with a bank of recorders forty kilometres out, and I was, as I have said, bored, because the bomb was the one event in that desert whose shape I could already predict. It would be enormous and it would be simple. Enormous and simple things do not reward attention.\par

There were nine of us in that tent, and eight of them were better scientists than I was, and every one of the eight spent the last minute of the count with his eyes fixed on the same three drums, the ones aimed at the tower, waiting to watch the desert's own hand thrown up across the paper. When it came, the light came first, through the canvas, a white with no morning anywhere in it; it threw the shadow of the tent-pole across the ground in the wrong direction for the smallest part of a second, and then took it back. Then the ground, a long low shove that lifted the recorders on their table and set them down again. Then, much later, absurdly, the sound, arriving like a thing that had forgotten itself and come running, a flat enormous slap laid against the whole horizon at once. The eight men watched the great strokes climb and break across the three drums, and were moved, I think, in the way men are moved by their own audacity. I kept my eyes on the other drums. The small ones. The ones aimed at the parts of the night where nothing was supposed to be.\par

So while the others watched the paper for the great sharp shout of the device, I did the thing I have done all my life and been called difficult for: I watched the quiet parts. The margins. The minutes before, when the desert was supposedly empty. The hours after, when it was supposedly still ringing down from the shot and nothing new should be able to be heard through the ring.\par

And in the hours after, while the great signal was still fading off the drums and every trained eye in the tent was following it down, there was something else on my papers.\par

It was small. It was so small that on any ordinary day it would have been thrown out as the tremor of a passing truck or the settling of the recorder's own frame. And it came from the wrong place: not from the tower, not from any fault the maps knew, but from far to the south, out of a stretch of the Tanezrouft where the maps put nothing but distance. And it had a shape I had never seen on a drum before, and I had by then seen a great many drums.\par

It was not the sharp single shout of an explosion. It was not the ragged rise and fall of a quake. It came in parts. Five of them. Five small clean strokes, evenly weighted, one after another, with a spacing between them that did not vary, and then it stopped, complete, the way a sentence stops. And beneath the five, or wrapped around them, there was a longer, softer return, a swelling and a dying, as though the five small blows had woken something very large and very far away and it had rolled over once in its sleep and settled.\par

I have carried that swelling my whole life, and for the whole of it I read only the half of what it was. I took the five even strokes for the event, and the long soft roll behind them for the ground settling under the blow --- an aftershock, a nothing, the kind of thing a careful reader is trained to set aside. That the roll was the greater half of it, and the five only the knock that woke it, I did not let myself understand for sixty-six years; and when at last I did, these last days, in this house, it was not off any drum of mine that I read it. But I run ahead of myself. That first night I understood only one thing, and it was enough to spend the rest of my life on. Nothing in nature does a thing in five even parts and then stops. Nature does not count. Nature has no fives. A thing that arrives in a fixed number of even strokes, at a fixed spacing, from a place that holds nothing, has been made to do that by a hand.\par

I did not tell anyone. I want that understood, because it is the beginning of what I am, and I will not have it dressed up later as scruple. I did not tell anyone because it was mine. In the space of one drum-length I had found the only genuinely new thing in a desert full of men making the loudest old thing there is, and I looked at it, and my first clear feeling --- before wonder, before fear, before any decent thing --- was that no one else in that tent was going to be allowed to see it. I folded the paper along the margin, and I dated it, and I put it inside my coat, and I sat through the rest of that historic night with the greatest secret of the century against my ribs and a face arranged for a woman watching a bomb.\par

\begin{room}

--- Forgive me. I am only a fact-checker, and it is my trade to ask the small stupid question. A great camp had come into that desert. Trucks, generators, drilling plant, machines by the hundred, run by men working through the night. Could your five strokes not have been some engine of their own, that no one had thought to write down?

--- It is the right question, and I put it to myself for four years before I would let the answer stand. An engine runs. It does not strike five even blows and then keep silence for four years. A generator hums; it does not compose a figure and sign it and put down the pen. And their camp was to the north, at the tower. My small thing came out of the south, out of a quarter where the French had set nothing at all, because there was nothing there they knew to set it on. It was not their machine. It was a machine that had been running in that ground long before theirs arrived, and would be running long after they left, and it had only been caught because they lit an enormous light beside it, and I was the one person in the desert looking the wrong way.

\end{room}

\ornament

That was one event, and one event is nothing. One event is a fault of the instrument, a fluke of the rock, a truck you did not know about. I knew that better than anyone, because knowing it is the whole of the discipline. A single wonder is only a mistake you have not yet found. So I did the correct thing, which is also the cruel thing: I waited. I let it be one event, and I built the rest of my life into a shape that would still be standing there, ear turned to that empty quarter, whenever the second one came.\par

It came four years later, when the French had moved their testing into the mountains of the Hoggar and buried it underground, and the whole apparatus of listening had moved and multiplied with it, and I had married my access into a permanent place inside it. Another shot, another night of great signals, another crowd of trained men following the loud thing down. And in the quiet after, from the south, from the same empty quarter, to the width of a hair the same: five even strokes, the same spacing, the same long soft roll waking behind them. I laid the new paper over the old one I had carried for four years, folded and refolded until it had gone soft, and the two drawings lay on each other the way a face lies on its reflection.\par

I did not celebrate. I remember precisely what I did: I locked the door of the room I was in, and I was sick, quietly and thoroughly, into a metal bin, and then I washed my face and went back to work. What turned my stomach was the size of the thing. I had spent four years telling myself that the likeliest truth was that I had seen nothing at all, a smudge, a fault, the one error every young reader makes and spends a career being humbled by. And the paper in my two hands had just taken that comfort away, and left me alone with the other possibility, the one I had wanted, God help me, more than I had ever wanted anything, and which was going to cost me every ordinary thing a life is supposed to be furnished with. I was right. That was the whole of it. I was right, and I have not once, in the sixty-six years since, been let stop being right.\par

A thing that happens once is an accident. A thing that happens twice, identically, is a machine. There is no third possibility, and I did not need a third to know what I had, but the world was generous and gave me a third some years on, and a fourth, at long, patient, maddening intervals --- always after the loud events, always small, always from the same nothing, always the same five strokes and the same great sleeping roll. The machine ran slowly. It struck its five notes and then went silent for years. Whatever it was, out there, it did not want to be found, and it had achieved that not by hiding its voice but by using it so rarely, and so quietly, that only a person who had already decided it was there, and had built a life out of waiting, would ever gather enough of it to be sure. I was that person. I had made myself that person over a bomb, at twenty-seven, out of boredom and appetite, and I do not repent of it, and you would not have wanted me to, because the incurious version of me would have folded that first paper into a report, and it would have gone up a chain of men who would have done to your Mountain, in one loud impatient decade, what I have spent sixty-six careful years preventing.\par

\begin{room}

--- She wants you to be grateful. Do you hear it? She has just told us she hunted my dead for sixty-six years, and she wants the room to thank her for being unhurried about it.

--- I hear it.

--- And you write it down anyway.

--- I write all of it down. That is the one thing left in this house that is still honest.

\end{room}

The strokes told me a machine existed and roughly where. They could tell me nothing of what it was, or who ran it, and I could not simply walk into an empty quarter of the Sahara and dig, because the empty quarter was a great deal larger than my patience, and because a woman with a fortune and a set of seismic papers announcing that there was a secret installation in the deep desert would have had the papers taken from her within a season by one of several governments and never seen again, and I would have become a footnote in someone else's file. So I did what my father would have done, and went to the cellar.\par

Edward's records were in three languages and no order, forty years of a man moving wealth by handshake and threat and midnight telephone, and I had them shipped to me and I read them the way I read the drums, for the quiet parts. And the quiet parts of my father's business, the failures, the dead ground, the money that went south and came back as nothing, all of it, when I laid it on a map, drew a ring. A ring of avoidance. Every route bent around one place. Every survey that touched a certain reach of the Tanezrouft returned corrupted or did not return. A drilling crew that reached a certain dry wadi came home a week early, every man of them, and gave notice inside the month, and would not be hired again for that country at any wage. A geologist of my father's, a sober Scot who had walked the Rub al-Khali and feared nothing that came out of sand, put in a report that a certain sixty kilometres of the Tanezrouft were, in his professional judgement, to be surveyed from the air and never once on foot; he gave no reason for it, and drank himself out of my father's service rather than be made to give one. Money went south and became nothing, again and again, across fifteen years; and my father, who could tell you to the franc where every other coin he had spent in his life had gone, could not account for that money. It was the single hole in the ledger of the one man I have known who did not permit himself holes. Two guides had refused a certain line at a certain latitude for wages that would have bought their villages. And in the margins, in my father's own furious hand, beside three separate failures across fifteen years, one word, underlined the third time hard enough to tear the paper. A name. Not a place. A person.\par

Amastan.\par

My father had circled the same emptiness I had, from the other side, with money instead of instruments, and he had come to the same edge and been turned back, and the thing that turned him back was not the desert. It was a man. The desert does not refuse wages. Men refuse wages. Somewhere between my father's cellar and my drum-papers there was a human door, and it had a name, and the name was the first warm thing in the whole cold problem, and I went toward it the way you go toward the one lit window in a dark street.\par

\begin{room}

--- Amastan. That is a name of the veiled people. A man of the sand. \ldots{} You are telling me the Mountain kept a hand out in the desert the whole time. Men not of us, who knew enough to turn a rich man back three times, and stayed poor enough to be trusted with the knowing. A ring beyond the ring. Shield would not have told us. A Veil is let to believe the desert is the last wall. If we knew there were hands past the wall, we would ask whose, and how many, and how bought, and Shield does not want a Veil asking who can be bought.

--- Your father asked exactly those questions, and the asking is what turned him back. I asked a smaller one. I have never once been frightened by the existence of a door. Only by the chance that another hand might reach it before mine.

\end{room}

My husband did not come with me the rest of the way, and I should tell you why, since you are keeping the record, and since it is the part I am least admired for.\par

He gave me everything I have described and one thing more, which was the door into the French machinery of secrets. A man who makes a nation its bomb is owed favours by the men who guard the nation, and he spent some of those favours, without asking closely why, on his clever wife who wanted access to seismic archives and desert logistics and, later, to certain quiet arrangements with certain quiet men. He opened those doors for me because I let him believe we were walking through them together. We were not. I had understood, by then, the one law your Mountain and I share, though I did not yet know it was your Mountain's law: that a thing known by two is a thing half-lost, and a thing known by one is a thing kept. Everything he gave me I received. Nothing I found did I give back. He thought we were building a house with two doors. I was building a house with one, and putting him outside it a room at a time, so gently and so reasonably that by the end he was grateful for the fresh air.\par

He died before any of it came to its harvest. The machinery that makes suns is not gentle with the men who serve it closest, and it was not gentle with him; he had stood too near too many of his own answers, and the science that was being born to read the ground had a sister science, younger and less spoken of, that was learning to read what such standing does to a man's blood. He went in under a year, badly. I was with him for it. I want that in the record too, beside the rest: that I sat with the man I had hollowed out, and held what was left of his hand, and grieved him honestly, and that neither the hollowing nor the grief cancelled the other. Both are true. I have made my peace with rooms that hold two true things. It is the only religion I have ever been able to keep.\par

What he left me was his access, which did not lapse for years because the men who owed him were slow to notice he was dead, and one unfinished sentence he had said to me near the end, drugged and wandering, which I have never repeated and will not repeat now, except to say that it was about chains, and about how the smallest possible push, delivered in exactly the right place at exactly the right moment, need supply almost none of the energy of the thing it starts. He had spent his life on that sentence pointed at the atom. He did not live to hear that it was also the truest thing anyone had ever said about the five little strokes from the south. I finished the sentence for him. I have been finishing it for sixty-six years.\par

\ornament

The man who carried the French half of it after my husband was named Martinez. He is old now, older than I am in the ways that matter, and you will not meet him, so I will be brief and unkind, which is the only register he has ever respected. Roger Martinez came up under my father in the trade and stayed on to serve me, and what he served, under both of us, was never quite Blackwood and never quite France but the narrow overlapping shadow the two throw when they stand close, which is where the yellow rock is always found and where the paper is always shabby and where a certain kind of man is always needed to go himself and sign it. Wherever the deep desert had a secret worth a government's while, and a mineral worth a company's, there was Martinez, with the right introductions and the wrong conscience, holding a door open for whichever of his two masters arrived first. He gave me protection, and classification, and continuity, and the quiet understanding that a discovery of the kind I was making does not belong to the person who makes it but to the arrangement that lets her survive making it. I took his help for thirty years and I never once mistook it for loyalty, and he never once offered it as such, and we understood each other perfectly, and I have outlived my usefulness to his masters and he knows it and I know it, which is a thing I will come back to, because it is the reason this house is not as safe as this boy believes it to be.\par

But that comes later. On the day I am telling you of, Martinez was only the man who could find me a person in a desert. And I gave him a name out of my dead father's margin, and told him to bring me the man who had turned Edward Blackwood back three times: alive, and unfrightened, and above all unbought by anyone else first. His name was Amastan. And with the finding of that name the science ended, because everything after it is not seismology. It is only what a person will do to reach a thing they have decided is theirs.\par

\begin{room}

--- Notice, before we go on, what the first thing was that I did with the greatest discovery of the age. I took a name out of my dead father's margin, and I hired a man to help me commit the first of a great many crimes against your people. I am not telling you this to be forgiven; I gave up forgiveness a long way back, and I have not once missed it. I am telling you so that when I am finished, and you are the one deciding what is to be done with me, you will not have the excuse of having thought me better than I am.

--- \ldots{}

\end{room}
